% Bibliothèques SI simplifiées, compatibles avec les interros
\providecommand{\og}{``}
\providecommand{\fg}{''}

% Graphe de liaisons
\tikzset{
  glsolide/.style={circle,draw=SIblue,very thick,minimum size=10mm,fill=SIlightblue,font=\bfseries},
  glliaison/.style={draw=SIred,very thick,rounded corners=2pt,font=\small,align=center}
}
\newenvironment{grapheliaisons}[1][]{\begin{center}\begin{tikzpicture}[node distance=25mm,#1]}{\end{tikzpicture}\end{center}}
\newcommand{\GLsolide}[3]{\node[glsolide] (#1) at (#3) {#2};}
\newcommand{\GLliaison}[4][]{\draw[glliaison,#1] (#2) -- node[midway,fill=white,inner sep=2pt]{#4} (#3);}
\newcommand{\GLpivot}[4][]{\GLliaison[#1]{#2}{#3}{Pivot\\$#4$}}
\newcommand{\GLglissiere}[4][]{\GLliaison[#1]{#2}{#3}{Glissière\\$#4$}}
\newcommand{\GLencastrement}[3][]{\GLliaison[#1]{#2}{#3}{Encastrement}}
\newcommand{\GLhelicoidale}[4][]{\GLliaison[#1]{#2}{#3}{Hélicoïdale\\$#4$}}

% Schémas cinématiques simplifiés
\newenvironment{schemacinematique}[1][]{\begin{center}\begin{tikzpicture}[line cap=round,line join=round,#1]}{\end{tikzpicture}\end{center}}
\newcommand{\SCbati}[2]{\coordinate (#1) at (#2);\draw[thick] ($(#1)+(-.65,0)$)--($(#1)+(.65,0)$);\foreach \x in {-.55,-.3,...,.55}{\draw ($(#1)+(\x,0)$)--++(-.18,-.18);}}
\newcommand{\SCpivot}[2]{\node[draw,circle,minimum size=5mm,thick] (#1) at (#2) {};}
\newcommand{\SCglissiere}[2]{\node[draw,rectangle,minimum width=10mm,minimum height=5mm,thick] (#1) at (#2) {};\draw[thick] ($(#1)+(-.35,-.18)$)--($(#1)+(.35,.18)$);}
\newcommand{\SChelicoidale}[2]{\node[draw,circle,minimum size=5mm,thick] (#1) at (#2) {};\draw[thick] ($(#1)+(-.15,0)$) arc (180:-90:.15);}
\newcommand{\SCrelier}[3][]{\draw[very thick,#1] (#2)--(#3);}

% SysML simplifié
\newenvironment{diagrammeSysML}[1][]{\begin{center}\begin{tikzpicture}[node distance=18mm,#1]}{\end{tikzpicture}\end{center}}
\newcommand{\SYSacteur}[4][]{\node[draw,circle,minimum size=3mm,#1] (head#4) at (#2) {};\draw (head#4)--++(0,-.55) coordinate (body#4);\draw ($(body#4)+(0,.35)$)--++(-.28,-.22);\draw ($(body#4)+(0,.35)$)--++(.28,-.22);\draw (body#4)--++(-.25,-.35);\draw (body#4)--++(.25,-.35);\node[below=8mm of head#4] (#4) {#3};}
\newcommand{\SYScas}[4][]{\node[draw,ellipse,align=center,minimum width=28mm,minimum height=10mm,#1] (#4) at (#2) {#3};}
\newcommand{\SYSinclude}[2]{\draw[dashed,->] (#1)--node[above]{\scriptsize\og include\fg}(#2);}
\newcommand{\SYSextend}[2]{\draw[dashed,->] (#1)--node[above]{\scriptsize\og extend\fg}(#2);}

% Tableau de mobilités
\newcommand{\mobilites}[6]{\begin{tabular}{c|c}R&T\\\hline #1&#2\\#3&#4\\#5&#6\end{tabular}}
